\subsection{A proper model for the higher-rank fibres}
\label{nf:graph-setup}

We retain $W=A\otimes B$, $S=\Seg(\PP A\times\PP B)$ and the
span variety $H\subseteq\Gr(r,W)$. The six matrices $Q_i$, their
five-dimensional span $L_0$, and $v(z)=(1,z,z^2)^{\mathsf T}$ are as
above. Put $T_5=\sum_{i=1}^5Q_i\otimes f_i$. In all the formats below,
\[
 a,b\geq3,\qquad a+b\geq r+2,\qquad c\geq r.
\]
The unordered span map is birational in this range. Indeed, choose
factor lists such that every subset has the largest possible rank.
Such lists exist by taking Vandermonde vectors. A linear combination
using exactly $s\geq2$ terms then has matrix rank at least
\[
 \min(a,s)+\min(b,s)-s\geq2,
 \qquad 2\leq s\leq r,
\]
by~\eqref{geom:sylvester}. The $r$ matrices are independent and their
span contains precisely their own rank-one directions. Thus a general
unordered tuple is uniquely recovered from its span. In characteristic
zero this gives birationality, as in the proof of
Proposition~\ref{prop:geometry-summary}.

We shall use a proper model which also records dependent limiting
configurations. Let $\mathcal K_r$ be the reduced closure of the graph
of the span map on independent reduced length-$r$ subschemes of $S$ in
\[
 \operatorname{Hilb}^r(S)\times H.
\]
The projection $\pi:\mathcal K_r\to H$ is proper and birational.
Properness follows from projectivity of $\operatorname{Hilb}^r(S)$;
birationality follows from that of the unordered span map.

\begin{lemma}[Fibre schemes and the normalization model]
\label{nf:proper-graph}
Suppose that $L\in H$ and the scheme $S\cap\PP L$ consists of
$m\geq r$ distinct reduced points. Then the fibre $\pi^{-1}(L)$ is a
reduced finite scheme, each of whose points is an $r$-element subset of
$S\cap\PP L$. Every independent such subset occurs. If
$\mathcal K_r$ is smooth at every point of this fibre, then, after
shrinking $H$ about $L$, the morphism $\pi$ is its normalization and
is unramified, with smooth source. The corresponding conclusions and
fibre cardinality hold over every rank-$r$ tensor with flattening image
$L$.
\end{lemma}

\begin{proof}
The condition that a finite subscheme $Z\subseteq S$ be contained in
a varying linear space $\PP K$ is closed. It holds on the ordinary
span graph and hence on $\mathcal K_r$. Therefore $\pi^{-1}(L)$ is
a closed subscheme of
\[
 \operatorname{Hilb}^r(S\cap\PP L).
\]
For a reduced set of $m$ points this Hilbert scheme is the disjoint
union of $\binom mr$ copies of $\Spec\C$, indexed by its
$r$-element subsets. Every closed subscheme of it is reduced. This
proves the fibre assertion, including the absence of nonreduced
collisions and of fibre multiplicities greater than one. Independent
subsets belong to the ordinary graph.

Since the fibre is finite, properness permits shrinking the target
until $\pi$ is finite. If the source is smooth along that fibre,
remove the proper image of its closed nonsmooth locus as well. The
resulting finite birational morphism has normal source and is thus the
normalization: its affine coordinate rings are integral over those of
the target, have the same fraction field, and are integrally closed.
At each point over $L$ the fibre cotangent space is zero. Base change
for relative differentials and Nakayama's lemma show that $\pi$ is
unramified there. Removing the finite image of the complement of the
unramified locus gives the asserted neighborhood.

Finally, Lemma~\ref{geom:flattening-bundle} identifies the rank-$r$
flattening locus of $X_r(a,b,c)$ with the bundle of surjections
$C^*\to K$ over $H$. Locally on $H$ this is the product with the
open set of full-rank $r$-by-$c$ matrices. Its pullback to the finite
smooth normalization is finite, birational and smooth, hence normal;
it is therefore the normalization of the tensor neighborhood.
Unramifiedness is preserved by this base change, and the fibre over a
fixed surjection is the same reduced fibre as over its image $L$.
This is the tensor transfer in Proposition~\ref{norm:criterion}.
\end{proof}

Near a configuration of distinct points, ordering the points gives a
finite \'etale cover by an open subset of $S^r$. Thus smoothness of
$\mathcal K_r$ may be checked on its ordered graph. The permutation
action is free there, including at dependent configurations.

\subsection{Diagonal extensions: six points in every rank}
\label{nf:diagonal-subsection}

Choose complementary coordinate vectors $u_1,\ldots,u_t$ in $A$
and $w_1,\ldots,w_t$ in $B$, where $t=r-5$, and set
\begin{equation}\label{nf:diagonal-data}
 D_s=u_s\otimes w_s,\qquad
 L_{\mathrm{diag}}=L_0\oplus\langle D_1,\ldots,D_t\rangle,
 \qquad
 T_{\mathrm{diag}}=T_5+\sum_{s=1}^tD_s\otimes f_{5+s}.
\end{equation}

\begin{theorem}[Diagonal fibres]
\label{nf:diagonal}
For every $r\geq6$, $a,b\geq r-2$ and $c\geq r$, the tensor
$T_{\mathrm{diag}}$ has rank and border rank $r$. The normalization
of $X_r(a,b,c)$ has exactly six reduced points over it. All six are
smooth and correspond bijectively to its six actual unordered
$r$-term decompositions. There is a Zariski open neighborhood of
$T_{\mathrm{diag}}$ whose entire normalization is smooth and whose
normalization map is unramified. In particular, no boundary
normalization point occurs over this tensor. These conclusions hold
after any permutation of the three factors.
\end{theorem}

The key exclusion uses the following block calculation. It applies to
arbitrary rectangular outer blocks, and so includes all the permitted
padding of the factors.

\begin{lemma}[Diagonal block valuation]
\label{nf:diagonal-block}
Let $R=\C[[z]]$, and let $P,N$ be matrices over $R$ of rank at most
three over $\C((z))$, both specializing to $M$ in the upper-left
$3$-by-$3$ block and to zero elsewhere. Write
\[
 P=\begin{pmatrix}A_P&B_P\\ C_P&E_P\end{pmatrix},
 \qquad
 N=\begin{pmatrix}A_N&B_N\\ C_N&E_N\end{pmatrix}.
\]
If $A_P-A_N$, $B_P-B_N$ and $C_P-C_N$ are divisible by $z^k$,
where $k\geq1$, then $E_P-E_N$ is divisible by $z^{k+1}$.
\end{lemma}

\begin{proof}
The $A$ blocks are invertible over $R$. The rank bounds give the
exact identities
\[
 E_P=C_PA_P^{-1}B_P,\qquad E_N=C_NA_N^{-1}B_N.
\]
Both cross blocks of each matrix are divisible by $z$. Moreover,
\begin{align*}
 E_P-E_N
 &=C_PA_P^{-1}(B_P-B_N)
   +(C_P-C_N)A_N^{-1}B_N\\
 &\hspace{12mm}+C_P(A_P^{-1}-A_N^{-1})B_N,\\
 A_P^{-1}-A_N^{-1}
 &=-A_P^{-1}(A_P-A_N)A_N^{-1}.
\end{align*}
The first two summands are divisible by $z^{k+1}$ and the last by
$z^{k+2}$, proving the claim.
\end{proof}

\begin{proof}[Proof of Theorem~\ref{nf:diagonal}]
A matrix in $L_{\mathrm{diag}}$ is block diagonal, with one core
$3$-by-$3$ block and $t$ scalar diagonal blocks. The rank of such a
matrix is the sum of the ranks of its nonzero blocks. A nonzero
rank-one matrix consequently has only one nonzero block. By
Lemma~\ref{norm:core}, the projective rank-one intersection is
exactly
\begin{equation}\label{nf:diagonal-intersection}
 S\cap\PP L_{\mathrm{diag}}
 =\{[Q_1],\ldots,[Q_6],[D_1],\ldots,[D_t]\}.
\end{equation}
This equality holds scheme-theoretically with all points reduced.
Indeed, at a core point choose a nonzero core pivot and normalize both
factor coordinates at that pivot to one. The off-block entries in
its row and column force every complementary factor coordinate to
vanish scheme-theoretically. The remaining intersection is the
reduced core intersection. At $[D_s]$ use its diagonal entry as pivot:
the vanishing of every other entry in that pivot row and column forces
all other factor coordinates to vanish. The same argument eliminates
unused coordinates under arbitrary padding.

The only relation among the $r+1$ displayed generators is
$\sum_{i=1}^6Q_i=0$. Thus their independent $r$-subsets are exactly
the six subsets omitting one core point. Each is a basis of
$L_{\mathrm{diag}}$. The third flattening of $T_{\mathrm{diag}}$
has rank $r$ and image $L_{\mathrm{diag}}$. An $r$-term
decomposition must give a rank-one basis of that image, and conversely
each such basis determines the third factors uniquely. It therefore
gives exactly one unordered decomposition into nonzero pure summands.
There are exactly six such decompositions, and the displayed sum and
the flattening lower bound give rank and border rank $r$.

We must exclude the remaining $r$-subsets from the proper graph.
By Lemma~\ref{nf:proper-graph}, they are the only possible boundary
configurations. Fix one omitting $D_s$; it contains all six core
points and every other $D_j$. Suppose it occurs in $\mathcal K_r$.
Choose a curve through this point with generic point in the ordinary
independent span graph, normalize the curve and complete its local
ring. After lifting to the ordering cover, this gives rank-one matrix
arcs over $R=\C[[z]]$,
\[
 p_i(z)\longrightarrow Q_i\quad(1\leq i\leq6),\qquad
 d_j(z)\longrightarrow D_j\quad(j\ne s),
\]
whose generic span has dimension $r$ and specializes to
$L_{\mathrm{diag}}$. Representatives regular in their pivot charts
may be chosen with exactly the indicated limits.

The first-factor vectors of the retained $d_j(z)$ are independent
over $R$, as are their second-factor vectors. Completing these lists
to bases gives changes of basis in $A$ and $B$, congruent to the
identity modulo $z$, which make all retained $d_j(z)$ equal to $D_j$.
These changes preserve the limiting subspace. The moving span is a
rank-$r$ subbundle of the constant $W$-bundle. Choose a regular section
$Z(z)$ of it satisfying $Z(0)=D_s$, and write over $\C((z))$
\[
 Z(z)=\sum_{i=1}^6\alpha_i(z)p_i(z)
                  +\sum_{j\ne s}\beta_j(z)D_j.
\]
Some coefficient has a pole, since the limiting tuple does not span
$D_s$. Let $k\geq1$ be the largest pole order. After multiplication
by $z^k$, reduction modulo $z$ gives a nonzero relation on the limiting
tuple. The six core coefficients therefore all have the same nonzero
constant $\lambda$, and every retained diagonal coefficient has
constant zero. Dividing by $\lambda$, denote the integral
coefficients by $a_i(z)$ and $b_j(z)$; then
\[
 a_i(0)=1,\qquad b_j(0)=0.
\]
Set
\[
 P=\sum_{i=1}^3a_ip_i,\qquad
 N=-\sum_{i=4}^6a_ip_i.
\]
Both matrices have rank at most three and specialize to $M$ in the
core block. Their relation is
\begin{equation}\label{nf:diagonal-pole-relation}
 P-N+\sum_{j\ne s}b_jD_j=\lambda^{-1}z^kZ(z).
\end{equation}
Its core and cross blocks show that the three differences in
Lemma~\ref{nf:diagonal-block} are divisible by $z^k$. The whole
lower-right block of $P-N$ is consequently divisible by $z^{k+1}$.
At the $(u_s,w_s)$ coordinate, all retained $D_j$ vanish, whereas
the right side of~\eqref{nf:diagonal-pole-relation} has valuation
exactly $k$, because $Z(0)=D_s$. This is a contradiction.

Thus the proper fibre contains only the six independent subsets.
At each, the graph is an open chart in the smooth unordered
configuration space. Lemma~\ref{nf:proper-graph} now proves the
normalization assertions, including the multiplicity one of every
fibre point and the tensor-neighborhood statement. The dimension
assumptions give $a+b\geq2r-4\geq r+2$, as required for generic
birationality. Permuting tensor factors gives the final assertion.
\end{proof}

\subsection{The seventh point at rank six}
\label{nf:rank-six-subsection}

\begin{theorem}[The rank-six row-extension fibre]
\label{nf:rank-six}
Let $a\geq4$, $b\geq3$, $a+b\geq8$ and $c\geq6$. For a new
coordinate vector $u\in A$ outside $A_0$, put
\[
 R=u\otimes v(4),\qquad L=L_0\oplus\C R,\qquad
 T_{\mathrm{row},6}=T_5+R\otimes f_6.
\]
The tensor has rank and border rank six. The normalization of
$X_6(a,b,c)$ has exactly seven reduced points over it, all smooth.
Six correspond to actual six-term decompositions. The seventh
corresponds to the dependent configuration of all six core directions
$[Q_1],\ldots,[Q_6]$ and is not an actual decomposition. On a
Zariski open neighborhood of this tensor the entire normalization is
smooth and its morphism is unramified. Exchanging $A$ and $B$ gives
the column-extension statement with $a\geq3$, $b\geq4$,
$a+b\geq8$ and $c\geq6$; all factor permutations are allowed.
\end{theorem}

\begin{proof}
Proposition~\ref{norm:row-column} gives the reduced intersection
\[
 S\cap\PP L=\{[Q_1],\ldots,[Q_6],[R]\},
\]
the six actual decompositions and the rank assertions. The six
subsets omitting a core point are independent. By
Lemma~\ref{nf:proper-graph}, the only remaining candidate in the
proper fibre is the six-core configuration. We calculate the graph
scheme at that configuration.

Let $\mathcal C_6\subseteq S^6$ be the ordered distinct-point
configuration space, of dimension $n=6(a+b-2)$, and let
$x=([Q_1],\ldots,[Q_6])$. Choose regular nonzero representatives
$P_i$ near $x$ and a fixed invertible five-row coordinate minor of
the first five columns of $(P_1,\ldots,P_6)$. Subtract from $P_6$
the unique linear combination of the first five columns which agrees
on those pivot rows, and call the residual vector $q$. Its remaining
coordinates generate the ideal $I$ of all six-by-six minors. This is
an equality of ideals, since row reduction using the invertible
pivot changes the column matrix to an identity block with final
residual column $q$. Write $D=V(I)$ for this exact determinantal
subscheme.

Equivalently, impose the relation
\begin{equation}\label{nf:rank-six-relation}
 \sum_{i=1}^6\lambda_iP_i=0,\qquad \lambda_6=1.
\end{equation}
The pivot solves the other five coefficients uniquely and regularly,
so eliminating them identifies this incidence scheme with $D$.
All six coefficients at $x$ are nonzero; we may restrict to the
locus where they are units. Lift the projective points to nonzero
factor vectors $a_i,b_i$. This is a smooth torus torsor of relative
dimension twelve. Absorbing $\lambda_i$ into $a_i$ turns
\eqref{nf:rank-six-relation} into
\begin{equation}\label{nf:rank-six-product}
 \mathsf A\mathsf B^{\mathsf T}=0,
 \qquad
 \mathsf A=(a_1,\ldots,a_6),\quad
 \mathsf B=(b_1,\ldots,b_6),
\end{equation}
with the five relation coordinates otherwise free.

Here is an explicit scheme chart for~\eqref{nf:rank-six-product}.
Write
\[
 \mathsf A=[A_1\ A_2],\qquad
 \mathsf B^{\mathsf T}=\begin{bmatrix}B_1\\B_2\end{bmatrix},
 \qquad B_1=[D_0\ E_0],\quad B_2=[C_0\ F_0],
\]
where $C_0,D_0$ are three-by-three. At the core, $A_1$ has a
fixed invertible three-row minor and $C_0$ is invertible. On that
open chart the product equations are exactly
\[
 A_1D_0+A_2C_0=0,\qquad A_1E_0+A_2F_0=0.
\]
They are scheme-theoretically equivalent to
\begin{equation}\label{nf:rank-six-solved-chart}
 A_2=-A_1D_0C_0^{-1},\qquad
 E_0=D_0C_0^{-1}F_0.
\end{equation}
Indeed, the first equation solves $A_2$; after substitution, a left
inverse for $A_1$ turns the second into the displayed equation for
$E_0$. There are no further equations. The free coordinates
$A_1,C_0,D_0,F_0$ have total dimension
\[
 3a+9+9+3(b-3)=3a+3b+9.
\]
Adding the five relation coordinates and descending through the
12-dimensional smooth torsor shows that $D$ is smooth of dimension
$3a+3b+2$. Its codimension in $\mathcal C_6$ is
\begin{equation}\label{nf:rank-six-center-codimension}
 h=3a+3b-14.
\end{equation}
In particular the exact maximal-minor ideal, and not merely its
radical, is the ideal of a smooth center.

At $x$, differentiating the residual column gives, up to a common
nonzero scalar,
\[
 dq_x(\delta P_1,\ldots,\delta P_6)
       =\sum_{i=1}^6\delta P_i\pmod{L_0}.
\]
The sum of the affine rank-one tangent spaces at the $Q_i$ is
\[
 U=A_0\otimes B+A\otimes B_0,
 \qquad \dim U=3a+3b-9.
\]
To see the equality, use the tangent formula
$T_{a_i\otimes b_i}\widehat S=a_i\otimes B+A\otimes b_i$,
and the facts that the $a_i$ span $A_0$ and the $b_i$ span $B_0$.
Projective representative changes discard only radial directions,
which lie in $L_0$. Thus the differential image is exactly $U/L_0$,
of dimension $h$. Its kernel is $T_xD$, and it induces an isomorphism
\begin{equation}\label{nf:rank-six-normal-map}
 N_{D/\mathcal C_6,x}\simeq U/L_0.
\end{equation}

The Pl\"ucker coordinates of the rational span map are
\[
 P_1\wedge\cdots\wedge P_5\wedge P_6
       =P_1\wedge\cdots\wedge P_5\wedge q.
\]
Their base ideal is exactly $I$. The closure of this rational graph
is therefore the blowup $\operatorname{Bl}_D(\mathcal C_6)$:
on affine charts the Rees algebra $\bigoplus_{j\geq0}I^j$ embeds
in the polynomial algebra with its degree parameter, and its Proj
is the closure of the graph where $I$ is nonzero. The Pl\"ucker
relations hold on this closure, so using the Grassmannian as target
gives the same graph. Both $\mathcal C_6$ and $D$ are smooth, hence
this blowup is smooth. Explicitly, in \'etale local coordinates with
$I=(x_1,\ldots,x_h)$, the chart indexed by $x_j$ has equations
$x_i=x_jz_i$ for $i\ne j$ and is smooth after eliminating those
$x_i$.

The exceptional fibre over $x$ is $\PP(N_{D/\mathcal C_6,x})$.
Under~\eqref{nf:rank-six-normal-map}, its map to the Grassmannian is
\[
 [w]\longmapsto L_0+\C w,\qquad [w]\in\PP(U/L_0).
\]
Different lines give different six-dimensional subspaces. Since
$R\in A\otimes B_0\subseteq U$ and $R\notin L_0$, exactly one
exceptional point maps to $L=L_0+\C R$. It is smooth. The free
permutation quotient preserves this local smoothness and gives
exactly one point over the unordered six-core configuration.

The complete proper fibre consequently consists of this point and
the six smooth ordinary points. By Lemma~\ref{nf:proper-graph},
all seven have multiplicity one, and they give the entire smooth,
unramified normalization over a neighborhood of $L$ and then over
a neighborhood of $T_{\mathrm{row},6}$. The exceptional configuration
spans only $L_0$, so it cannot represent an actual six-term
decomposition of a tensor whose third flattening has rank six.
The symmetry statements follow by permuting the factors.
\end{proof}

\subsection{All ranks with a three-dimensional matrix factor}
\label{nf:three-factor-subsection}

\begin{theorem}[The three-dimensional-factor row family]
\label{nf:three-factor}
Let $r\geq6$, $a\geq r-1$ and $c\geq r$, and let $B=B_0$ have
dimension exactly three. For independent extra coordinate vectors
$u_1,\ldots,u_{r-5}$ in $A$, set
\begin{equation}\label{nf:row-data}
 R_s=u_s\otimes v(s+3),\qquad
 L_r=L_0\oplus\langle R_1,\ldots,R_{r-5}\rangle,\qquad
 T_{\mathrm{row}}=T_5+\sum_{s=1}^{r-5}R_s\otimes f_{5+s}.
\end{equation}
The tensor has rank and border rank $r$. The normalization of
$X_r(a,3,c)$ has exactly $r+1$ reduced points over it, all smooth.
Six correspond to actual $r$-term decompositions. The other $r-5$
correspond to the dependent configurations which omit one $R_s$
and contain all six $Q_i$. There is a Zariski open neighborhood of
$T_{\mathrm{row}}$ whose entire normalization is smooth and whose
normalization morphism is unramified. All factor permutations of
this statement hold.
\end{theorem}

\begin{proof}
By Proposition~\ref{norm:row-column}, the rank-one intersection is
exactly the $r+1$ listed reduced points, and the only relation among
them is the six-core relation. There are six independent
$r$-subsets, giving precisely the six actual decompositions; the
rank and border rank are $r$. The other $r-5$ subsets each omit
one added point. Lemma~\ref{nf:proper-graph} reduces exhaustion and
smoothness to the graph at these dependent configurations.

Fix one such ordered configuration $x=(p_1,\ldots,p_r)$ and reorder
it so that the first $r-1$ points are independent. Write $K$ for
their $(r-1)$-dimensional span at $x$. In local projective charts
choose nonzero affine representatives $P_i(y)$ and invert a fixed
$(r-1)$-row minor of the first $r-1$ columns. Their spans form a
subbundle $K(y)$ of the constant $W$-bundle, whose quotient is locally
free of rank
\[
 h=\dim W-(r-1)=3a-r+1.
\]
In a quotient trivialization let $q(y)$ be the residual image of
$P_r(y)$. Elementary row operations using the pivot show that the
coordinates of $q$ generate the ideal of all $r$-by-$r$ minors.
Thus the rank-deficiency subscheme is exactly $D=\{q=0\}$,
with this ideal, not just this zero set.

In the unique relation $\sum_i\lambda_iP_i(x)=0$, the six core
coefficients are nonzero and those of the retained added points are
zero. Differentiating the residual column yields, up to a common
nonzero scalar and quotient coordinates,
\begin{equation}\label{nf:row-residual-derivative}
 dq_x(\delta P_1,\ldots,\delta P_r)
   =\sum_{i\,\text{core}}\lambda_i\delta P_i\pmod K.
\end{equation}
The three row-core matrices are $Q_i=e_i\otimes r_i$, where the
rows $r_1,r_2,r_3$ of the invertible matrix $M$ form a basis of $B$.
Their affine tangent spaces contain $A\otimes r_i$; their sum
therefore contains all of $A\otimes B=W$. Each relevant coefficient
in~\eqref{nf:row-residual-derivative} is nonzero. Passing to
projective tangent spaces removes only the radial directions $Q_i$,
which already lie in $K$. Consequently $dq_x$ is surjective onto
$W/K$.

It follows from the Jacobian criterion that $q$ is smooth near $x$
and that $D$ is smooth of codimension $h$. Equivalently, the
coordinates $q_1,\ldots,q_h$ form part of an \'etale local coordinate
system. As in the rank-six calculation, the local ordered span graph
is the blowup of the exact maximal-minor ideal $(q_1,\ldots,q_h)$.
In its $q_j$ chart the equations are
\[
 q_i=q_jz_i\quad(i\ne j).
\]
Eliminating these $q_i$ leaves $q_j$, the $z_i$, and the complementary
smooth coordinates free. Thus the graph is smooth at every point
above $x$.

The exceptional fibre over $x$ is $\PP(W/K)$. Its map to the
Grassmannian is the usual isomorphism from $\PP(W/K)$ to the
variety of $r$-dimensional subspaces containing $K$,
\[
 [w]\longmapsto K+\C w.
\]
Since $K\subset L_r$, the prescribed target $L_r$ selects exactly
one exceptional point. In particular that point exists in the graph,
is smooth, and is the only point over this dependent configuration
and $L_r$. The ordering cover is \'etale, so the same is true in the
unordered graph.

This gives one point for each of the $r-5$ dependent subsets and one
ordinary smooth point for each of the six independent subsets.
Lemma~\ref{nf:proper-graph} proves that this is the entire reduced
fibre, that it has $r+1$ points, and that the normalization is smooth
and unramified on a neighborhood. Its tensor transfer proves the
assertions over $T_{\mathrm{row}}$. Here $a+3\geq r+2$ supplies
the required generic birationality. A dependent subset spans only an
$(r-1)$-space and hence cannot give an actual decomposition of this
rank-$r$ flattening.
\end{proof}

The equality $B=B_0$ is a hypothesis of
Theorem~\ref{nf:three-factor}: it is precisely what makes the three
row-core tangent spaces cover $A\otimes B$ in the residual derivative.
The first factor and the third factor may be enlarged throughout the
stated ranges. Theorem~\ref{nf:rank-six} separately permits enlargement
of both matrix factors at rank six. The three theorems above do not address the mixed row/column family
in ranks $r\geq7$ with both matrix factors enlarged; specified such
families will be treated below. Their neighborhood statements imply
no global smoothness, Cohen--Macaulayness or rational singularities
of a normalization.
